% =============================================================================
%  capitulos/cap4_requisitos.tex
%  CAPÍTULO 4 — REQUISITOS E MODELAGEM DO SISTEMA
%
%  Objetivo: provar que o sistema foi planejado com rigor de engenharia.
%  Aqui você demonstra a diferença entre "programador" e "engenheiro".
%
%  Contém: Requisitos Funcionais (RF), Requisitos Não Funcionais (RNF)
%          e os principais Diagramas UML.
% =============================================================================

\chapter{Requisitos e Modelagem do Sistema}
\label{cap:requisitos}

TODO: Parágrafo introdutório do capítulo, explicando como os requisitos
foram levantados (técnica utilizada: entrevistas, questionários, análise
de sistemas similares, etc.) e o que será apresentado nas próximas seções.


% -----------------------------------------------------------------------------
%  4.1 REQUISITOS FUNCIONAIS (RF)
%  Descrevem O QUE o sistema faz — funcionalidades observáveis pelo usuário.
%  Convenção de numeração: RF01, RF02... ou RF-001, RF-002...
%  Prioridade: Alta / Média / Baixa  (ou MoSCoW: Must / Should / Could / Won't)
% -----------------------------------------------------------------------------
\section{Requisitos Funcionais}
\label{sec:requisitos-funcionais}

Os requisitos funcionais descrevem as funcionalidades que o sistema deve
oferecer aos seus usuários. A \cref{tab:rf} apresenta os requisitos
funcionais levantados durante a fase de análise.

\begin{table}[H]
  \caption{Requisitos Funcionais do Sistema}
  \label{tab:rf}
  \centering
  % Ajuste as colunas conforme necessário. X = coluna de largura flexível.
  \begin{tabularx}{\textwidth}{llXl}
    \toprule
    \textbf{ID}   & \textbf{Nome}        & \textbf{Descrição}
                  & \textbf{Prioridade}  \\
    \midrule
    RF01 & TODO: Nome curto & TODO: O sistema deve permitir que o
           usuário [ação]...                    & Alta   \\
    RF02 & TODO: Nome curto & TODO: O sistema deve [ação]...
                                                & Alta   \\
    RF03 & TODO: Nome curto & TODO: O sistema deve [ação]...
                                                & Média  \\
    RF04 & TODO: Nome curto & TODO: O sistema deve [ação]...
                                                & Média  \\
    RF05 & TODO: Nome curto & TODO: O sistema deve [ação]...
                                                & Baixa  \\
    % Adicione quantas linhas forem necessárias
    \bottomrule
  \end{tabularx}
  \fonte{Elaborado pelo autor.}
\end{table}


% -----------------------------------------------------------------------------
%  4.2 REQUISITOS NÃO FUNCIONAIS (RNF)
%  Descrevem COMO o sistema se comporta: performance, segurança,
%  usabilidade, disponibilidade, manutenibilidade, escalabilidade...
%  Sempre que possível, inclua uma métrica mensurável na coluna "Critério".
% -----------------------------------------------------------------------------
\section{Requisitos Não Funcionais}
\label{sec:requisitos-nao-funcionais}

Os requisitos não funcionais definem restrições de qualidade e
comportamento do sistema. A \cref{tab:rnf} apresenta os requisitos não
funcionais identificados.

\begin{table}[H]
  \caption{Requisitos Não Funcionais do Sistema}
  \label{tab:rnf}
  \centering
  \begin{tabularx}{\textwidth}{llXl}
    \toprule
    \textbf{ID}  & \textbf{Categoria}   & \textbf{Descrição}
                 & \textbf{Critério}    \\
    \midrule
    RNF01 & Performance   & TODO: O sistema deve responder às requisições
            da API...                     & Até X ms \\
    RNF02 & Segurança     & TODO: Senhas devem ser armazenadas com hash...
                                           & bcrypt/Argon2 \\
    RNF03 & Usabilidade   & TODO: A interface deve ser responsiva...
                                           & Mobile-first \\
    RNF04 & Disponibilidade & TODO: O sistema deve estar disponível...
                                           & 99\% uptime \\
    RNF05 & Manutenibilidade & TODO: O código deve ter cobertura mínima
             de testes...                  & $\geq$ 70\% \\
    % Adicione mais linhas conforme necessário
    \bottomrule
  \end{tabularx}
  \fonte{Elaborado pelo autor.}
\end{table}


% -----------------------------------------------------------------------------
%  4.3 MODELAGEM UML
%  Inclua os diagramas mais relevantes para o entendimento do sistema.
%  RECOMENDAÇÃO: use ferramentas externas (draw.io, Astah, PlantUML,
%  Lucidchart) e exporte como PNG ou PDF para incluir com \includegraphics.
%
%  Diagramas sugeridos (inclua os que fizerem sentido para o seu projeto):
%    - Diagrama de Casos de Uso      → visão geral do que o sistema faz
%    - Diagrama de Classes           → estrutura do código/domínio
%    - Diagrama de Sequência         → fluxo de uma operação específica
%    - Diagrama de Atividades        → fluxo de um processo de negócio
%    - Diagrama de Componentes       → arquitetura em alto nível
% -----------------------------------------------------------------------------
\section{Modelagem UML}
\label{sec:uml}

TODO: Parágrafo introdutório, dizendo quais diagramas foram elaborados
e qual ferramenta foi utilizada para criá-los.


% --- Diagrama de Casos de Uso ---
\subsection{Diagrama de Casos de Uso}
\label{subsec:casos-de-uso}

TODO: Descreva os atores (quem interage com o sistema) e os principais
casos de uso antes de mostrar o diagrama. Explique o que cada ator pode
fazer e quais são os casos de uso mais relevantes.

\begin{figure}[H]
  \caption{Diagrama de Casos de Uso}
  \label{fig:casos-de-uso}
  \centering
  % Substitua o \fbox pelo \includegraphics quando tiver a imagem pronta:
  % \includegraphics[width=0.85\textwidth]{imagens/diagrama_casos_uso.png}
  \fbox{\rule{0pt}{6cm}\rule{0.85\textwidth}{0pt}} % ← Placeholder, remova depois
  \fonte{Elaborado pelo autor com TODO: [nome da ferramenta].}
\end{figure}

TODO: Após o diagrama, descreva os casos de uso mais importantes,
complementando o que a imagem mostra.


% --- Diagrama de Classes ---
\subsection{Diagrama de Classes}
\label{subsec:diagrama-classes}

TODO: Explique as principais classes do sistema, seus atributos,
métodos e relacionamentos antes de mostrar o diagrama.

\begin{figure}[H]
  \caption{Diagrama de Classes}
  \label{fig:diagrama-classes}
  \centering
  % \includegraphics[width=\textwidth]{imagens/diagrama_classes.png}
  \fbox{\rule{0pt}{7cm}\rule{0.9\textwidth}{0pt}} % ← Placeholder, remova depois
  \fonte{Elaborado pelo autor.}
\end{figure}


% --- Diagramas adicionais (descomente e adapte conforme necessário) ---

% \subsection{Diagrama de Sequência}
% \label{subsec:diagrama-sequencia}
% TODO: Escolha um fluxo importante (ex.: autenticação, realização de pedido)
% e mostre a interação entre os componentes ao longo do tempo.
%
% \begin{figure}[H]
%   \caption{Diagrama de Sequência — TODO: Nome do Fluxo}
%   \label{fig:diagrama-sequencia}
%   \centering
%   % \includegraphics[width=\textwidth]{imagens/diagrama_sequencia.png}
%   \fbox{\rule{0pt}{6cm}\rule{0.9\textwidth}{0pt}}
%   \fonte{Elaborado pelo autor.}
% \end{figure}


% \subsection{Diagrama de Atividades}
% \label{subsec:diagrama-atividades}
%
% \begin{figure}[H]
%   \caption{Diagrama de Atividades — TODO: Nome do Processo}
%   \label{fig:diagrama-atividades}
%   \centering
%   % \includegraphics[width=0.7\textwidth]{imagens/diagrama_atividades.png}
%   \fbox{\rule{0pt}{8cm}\rule{0.6\textwidth}{0pt}}
%   \fonte{Elaborado pelo autor.}
% \end{figure}
